\begin{exercises}
  \recommended \item 
    求代数余子式。
    \begin{equation*}
      T=\begin{mat}[r]
           1  &0  &2  \\
          -1  &1  &3  \\
           0  &2  &-1
         \end{mat}
    \end{equation*}
    \begin{exparts*}
       \partsitem \( T_{2,3} \)
       \partsitem \( T_{3,2} \)
       \partsitem \( T_{1,3} \)
    \end{exparts*}
    \begin{answer}
      \begin{exparts*}
       \partsitem \( (-1)^{2+3}\begin{vmat}[r]
                            1  &0  \\
                            0  &2
                          \end{vmat}=-2  \)
       \partsitem \( (-1)^{3+2}\begin{vmat}[r]
                            1  &2  \\
                           -1  &3
                          \end{vmat}=-5  \)
       \partsitem \( (-1)^{4}\begin{vmat}[r]
                           -1  &1  \\
                            0  &2
                          \end{vmat}=-2  \)
      \end{exparts*}  
    \end{answer}
  \recommended \item 
    求该矩阵的伴随矩阵。
    \begin{equation*}
      T=\begin{mat}[r]
           1  &0  &2  \\
          -1  &1  &3  \\
           0  &2  &-1
         \end{mat}
    \end{equation*}
    \begin{answer}
      % sage: M = matrix(QQ, [[1,0,2], [-1,1,3], [0,2,-1]])
      % sage: M.adjoint()
      % [-7  4 -2]
      % [-1 -1 -5]
      % [-2 -2  1] 
      {\renewcommand{\arraystretch}{1.2}
      \begin{equation*} 
        \adj(T)=
        \begin{mat}
          \begin{vmatrix}
            1 &3 \\
            2 &-1
          \end{vmatrix}
          &-\begin{vmatrix}
            0 &2 \\ 
            2 &-1
          \end{vmatrix}
          &\begin{vmatrix}
             0 &2 \\
              1 &3
           \end{vmatrix}                  \\
          -\begin{vmatrix}
             -1 &3 \\
             0 &-1
           \end{vmatrix}
          &\begin{vmatrix}
            1 &2 \\
            0 &-1
           \end{vmatrix}
          &-\begin{vmatrix}
            1 &2 \\
            -1 &3
           \end{vmatrix}                  \\
          \begin{vmatrix}
            -1 &1 \\ 
            0 &2
          \end{vmatrix}
          &-\begin{vmatrix}
            1  &0 \\
            0 &2
           \end{vmatrix}
          &\begin{vmatrix}
            1 &0 \\
            -1 &1
           \end{vmatrix}                  
        \end{mat}
        =
        \begin{mat}
          -7 &4 &-2 \\
          -1  &-1 &-5 \\
          -2 &-2 &1 
        \end{mat}
      \end{equation*}}
    \end{answer}
  \item 这个行列式等于 $0$。请按第一行展开计算它。
    \begin{equation*}
      \begin{vmat}[r]
         1  &2  &3  \\
         4  &5  &6  \\
         7  &8  &9
      \end{vmat}
    \end{equation*}
    \begin{answer}
      \begin{equation*}
        \begin{vmat}[r]
          1  &2  &3  \\
          4  &5  &6  \\
          7  &8  &9
        \end{vmat}=
        1\cdot(+1)\cdot
          \begin{vmatrix}
            5 &6 \\
            8 &9
          \end{vmatrix}
        +2\cdot(-1)\cdot
          \begin{vmatrix}
            4 &6 \\
            7 &9
          \end{vmatrix}
        +3\cdot(+1)\cdot
          \begin{vmatrix}
            4 &5 \\
            7 &8
          \end{vmatrix}
      \end{equation*}
      $\nbyn{2}$ 矩阵的公式给出
      $1\cdot (-3)-2\cdot (-6)+3\cdot(-3)=0$。
    \end{answer}
  \recommended \item 
    按指定方式展开求行列式
    \begin{equation*}
      \begin{vmat}[r]
         3  &0  &1  \\
         1  &2  &2  \\
        -1  &3  &0
      \end{vmat}
    \end{equation*}
    \begin{exparts*}
      \partsitem 沿第一行展开
      \partsitem 沿第二行展开
      \partsitem 沿第三列展开。
    \end{exparts*}
    \begin{answer}
      \begin{exparts}
         \partsitem \( 3\cdot(+1)\begin{vmat}[r]
                          2  &2  \\
                          3  &0
                       \end{vmat}
                 +0\cdot (-1)\begin{vmat}[r]
                          1  &2  \\
                         -1  &0
                       \end{vmat}
                 +1\cdot (+1)\begin{vmat}[r]
                          1  &2  \\
                         -1  &3
                       \end{vmat} =-13 \)
         \partsitem \( 1\cdot (-1)\begin{vmat}[r]
                          0  &1  \\
                          3  &0
                       \end{vmat}
                 +2\cdot (+1)\begin{vmat}[r]
                          3  &1  \\
                         -1  &0
                       \end{vmat}
                 +2\cdot (-1)\begin{vmat}[r]
                          3  &0  \\
                         -1  &3
                       \end{vmat} =-13 \)
         \partsitem \( 1\cdot (+1)\begin{vmat}[r]
                          1  &2  \\
                         -1  &3
                       \end{vmat}
                 +2\cdot (-1)\begin{vmat}[r]
                          3  &0  \\
                         -1  &3
                       \end{vmat}
                 +0\cdot (+1)\begin{vmat}[r]
                          3  &0  \\
                          1  &2
                       \end{vmat} =-13 \)
       \end{exparts}  
     \end{answer}
  \item 
    求\nearbyexample{ex:ExpLaPlace}中矩阵的伴随矩阵。
    \begin{answer}
      这就是 $\adj(T)$。
      \begin{align*}
         \begin{mat}
           T_{1,1}  &T_{2,1}  &T_{3,1} \\
           T_{1,2}  &T_{2,2}  &T_{3,2} \\
           T_{1,3}  &T_{2,3}  &T_{3,3}   
         \end{mat}
         &=\begin{mat}
            \begin{vmat}
               5  &6  \\ 8  &9
            \end{vmat}
            &-\begin{vmat}
               2  &3  \\  8  &9
            \end{vmat}
            &+\begin{vmat}
               2  &3  \\  5  &6
            \end{vmat}        \\[2.5ex]
            -\begin{vmat}
               4  &6  \\  7  &9
            \end{vmat}
            &+\begin{vmat}
               1  &3  \\  7  &9
            \end{vmat}
            &-\begin{vmat}
               1  &3  \\  4  &6
            \end{vmat}        \\[2.5ex]
            +\begin{vmat}
               4  &5  \\  7  &8
            \end{vmat}
            &-\begin{vmat}
               1  &2  \\  7  &8
            \end{vmat}
            &+\begin{vmat}
               1  &2  \\  4  &5
            \end{vmat}
         \end{mat}                                   \\
         &=\begin{mat}[r]
          -3  &6   &-3 \\
           6  &-12 &6  \\
          -3  &6   &-3
         \end{mat}
      \end{align*}
    \end{answer}
  \recommended \item 
    求与下列每个矩阵伴随的矩阵。
    \begin{exparts*}
      \partsitem \( \begin{mat}[r]
                 2   &1  &4  \\
                -1   &0  &2  \\
                 1   &0  &1
               \end{mat}  \)
      \partsitem \( \begin{mat}[r]
                 3  &-1  \\
                 2  &4
               \end{mat}  \)
      \partsitem \( \begin{mat}[r]
                 1   &1  \\
                 5   &0
               \end{mat}  \)
      \partsitem \( \begin{mat}[r]
                 1   &4  &3  \\
                -1   &0  &3  \\
                 1   &8  &9
               \end{mat}  \)
    \end{exparts*}
    \begin{answer}
      \begin{exparts}
        \partsitem 这就是伴随矩阵。 
          \begin{align*}
          \begin{mat}
            T_{1,1}  &T_{2,1}  &T_{3,1} \\ 
            T_{1,2}  &T_{2,2}  &T_{3,2} \\ 
            T_{1,3}  &T_{2,3}  &T_{3,3} 
          \end{mat}
          &=\begin{mat}
            \begin{vmat}[r]
              0  &2  \\
              0  &1 
            \end{vmat}
            &-\begin{vmat}[r]
               1  &4  \\
               0  &1
            \end{vmat}
            &\begin{vmat}[r]
               1  &4  \\
               0  &2 
            \end{vmat}        \\[2.5ex]
            -\begin{vmat}[r]
              -1  &2  \\
               1  &1
            \end{vmat}
            &\begin{vmat}[r]
               2  &4  \\
               1  &1
            \end{vmat}
            &-\begin{vmat}[r]
               2  &4  \\ 
              -1  &2
            \end{vmat}         \\[2.5ex]
            \begin{vmat}[r]
              -1  &0  \\
               1  &0
            \end{vmat} 
            &-\begin{vmat}[r]
               2  &1  \\
               1  &0
            \end{vmat}
            &\begin{vmat}[r]
               2  &1  \\
              -1  &0
            \end{vmat}
          \end{mat}                         \\
          &=\begin{mat}[r]
             0  &-1  &2  \\
             3  &-2  &-8 \\
             0  &1   &1            
          \end{mat}
          \end{align*}
        \partsitem 子式是 $\nbyn{1}$ 的。
          $\begin{mat}
             T_{1,1}  &T_{2,1} \\
             T_{1,2}  &T_{2,2}  
          \end{mat}
          =
          \begin{mat}
            \begin{vmat}
              4              
            \end{vmat}
            &-\begin{vmat}
              -1
            \end{vmat}        \\[1.5ex]
            -\begin{vmat}
               2              
            \end{vmat}
            &\begin{vmat}
               3
            \end{vmat}
          \end{mat}
          =
          \begin{mat}[r]
            4  &1  \\
           -2  &3
          \end{mat}$
        \partsitem 
          $\begin{mat}[r]
               0  &-1 \\
              -5  &1
          \end{mat}$
        \partsitem 子式是 $\nbyn{2}$ 的。
          \begin{align*}
          \begin{mat}
            T_{1,1}  &T_{2,1}  &T_{3,1} \\ 
            T_{1,2}  &T_{2,2}  &T_{3,2} \\ 
            T_{1,3}  &T_{2,3}  &T_{3,3} 
          \end{mat}
          &=\begin{mat}
            \begin{vmat}[r]
              0  &3  \\
              8  &9
            \end{vmat}
            &-\begin{vmat}[r]
              4  &3  \\
              8  &9
            \end{vmat}
            &\begin{vmat}[r]
              4  &3  \\
              0  &3 
            \end{vmat}       \\[2.5ex]
            -\begin{vmat}[r]
              -1  &3  \\
               1  &9
            \end{vmat}
            &\begin{vmat}[r]
               1  &3  \\
               1  &9              
            \end{vmat}
            &-\begin{vmat}[r]
               1  &3 \\
              -1  &3              
            \end{vmat}        \\[2.5ex]
            \begin{vmat}[r]
               -1  &0  \\
                1  &8              
            \end{vmat}
            &-\begin{vmat}[r]
                1  &4 \\
                1  &8              
            \end{vmat}
            &\begin{vmat}[r]
                1  &4  \\
               -1  &0              
            \end{vmat}
          \end{mat}                       \\
          &=
          \begin{mat}[r]
            -24  &-12  &12  \\
             12  &6    &-6   \\
             -8  &-4   &4
          \end{mat}
        \end{align*}
      \end{exparts}
    \end{answer}
\recommended \item
    用\nearbytheorem{th:MatTimesAdjEqDiagDets}求前一题中每个矩阵的逆矩阵。
    \begin{answer}
      \begin{exparts}
        \partsitem $(1/3)\cdot 
           \begin{mat}[r]
             0  &-1  &2  \\
             3  &-2  &-8 \\
             0  &1   &1            
          \end{mat}          
          =
          \begin{mat}[r]
            0  &-1/3  &2/3  \\
            1  &-2/3  &-8/3 \\ 
            0  &1/3   &1/3
          \end{mat}$
        \partsitem $(1/14)\cdot 
          \begin{mat}[r]
            4  &1  \\
           -2  &3
          \end{mat}
          =
          \begin{mat}[r]
            2/7  &1/14  \\
           -1/7  &3/14
          \end{mat}$
        \partsitem $(1/-5)\cdot 
          \begin{mat}[r]
               0  &-1 \\
              -5  &1
          \end{mat}
          =
          \begin{mat}[r]
            0  &1/5 \\
            1  &-1/5
          \end{mat}$
        \partsitem 该矩阵的行列式为零，因此没有逆矩阵。
      \end{exparts}
    \end{answer}
  \item 
    求与下面这个矩阵伴随的矩阵。
    \begin{equation*}
      \begin{mat}[r]
        2  &1  &0  &0  \\
        1  &2  &1  &0  \\
        0  &1  &2  &1  \\
        0  &0  &1  &2
      \end{mat}
    \end{equation*}
    \begin{answer}
      $\begin{mat}
        T_{1,1}  &T_{2,1}  &T_{3,1}  &T_{4,1}  \\ 
        T_{1,2}  &T_{2,2}  &T_{3,2}  &T_{4,2}  \\ 
        T_{1,3}  &T_{2,3}  &T_{3,3}  &T_{4,3}  \\ 
        T_{1,4}  &T_{2,4}  &T_{3,4}  &T_{4,4}  
      \end{mat}
      =
      \begin{mat}[r]
        4  &-3  &2  &-1  \\
       -3  &6   &-4 &2   \\
        2  &-4  &6  &-3  \\
        -1 &2  &-3  &4
      \end{mat}$
    \end{answer}
  \recommended \item
    沿第一行展开，推导 \( \nbyn{2}  \) 矩阵的行列式公式。
    \begin{answer}
       行列式
       \begin{equation*}
         \begin{vmat}
           a  &b  \\
           c  &d
         \end{vmat}
       \end{equation*} 
       按第一行展开给出    
       \( a\cdot (+1)\deter{d}+b\cdot (-1)\deter{c}=ad-bc \)
       （注意那两个 $\nbyn{1}$ 子式）。
     \end{answer}
  \recommended \item
    沿第一行展开，推导 \( \nbyn{3} \) 矩阵的行列式公式。
    \begin{answer}
       矩阵
       \begin{equation*}
         \begin{mat}
           a  &b  &c  \\
           d  &e  &f  \\
           g  &h  &i
         \end{mat}
       \end{equation*}
       的行列式如下。
       \begin{equation*}
         a\cdot \begin{vmat}
             e  &f  \\
             h  &i
          \end{vmat}
         -b\cdot \begin{vmat}
             d  &f  \\
             g  &i
          \end{vmat}
         +c\cdot \begin{vmat}
             d  &e  \\
             g  &h
          \end{vmat}
         =a(ei-fh)-b(di-fg)+c(dh-eg)
       \end{equation*}  
    \end{answer}
  \recommended \item 
   \begin{exparts}
     \partsitem 给出 \( \nbyn{2} \) 矩阵的伴随矩阵的公式。
     \partsitem 利用它推导逆矩阵的公式。
   \end{exparts} 
    \begin{answer}
      \begin{exparts}
        \partsitem
          $\begin{mat}
             T_{1,1}  &T_{2,1}  \\
             T_{1,2}  &T_{2,2}
            \end{mat}
            =\begin{mat}
            \begin{vmat}
              t_{2,2}
            \end{vmat}
           &-\begin{vmat}
               t_{1,2}
             \end{vmat}       \\
           -\begin{vmat}
               t_{2,1}
            \end{vmat}
           &\begin{vmat}
              t_{1,1} 
            \end{vmat}
          \end{mat}
          =\begin{mat}
             t_{2,2}  &-t_{1,2}  \\
            -t_{2,1} &t_{1,1}
           \end{mat}$
        \partsitem
          $(1/t_{1,1}t_{2,2}-t_{1,2}t_{2,1})\cdot
          \begin{mat}
             t_{2,2}  &-t_{1,2}  \\
            -t_{2,1} &t_{1,1}
           \end{mat}$
      \end{exparts}
    \end{answer}
  \recommended \item
    我们能沿对角线展开来计算行列式吗？
    \begin{answer}
      不能。
      下面这个行列式的值
      \begin{equation*}
        \begin{vmat}[r]
          1  &0  &0  \\
          0  &1  &0  \\
          0  &0  &1
        \end{vmat}=1
      \end{equation*}
      并不等于沿对角线展开的结果。
      \begin{equation*}
        1\cdot (+1)\begin{vmat}[r]
               1  &0  \\
               0  &1
             \end{vmat}
       +1\cdot (+1)\begin{vmat}[r]
               1  &0  \\
               0  &1
             \end{vmat}
       +1\cdot (+1)\begin{vmat}[r]
               1  &0  \\
               0  &1
             \end{vmat}=3
      \end{equation*}  
    \end{answer}
  \item 
    给出对角矩阵的伴随矩阵的公式。
    \begin{answer}
      考虑这个对角矩阵。
      \begin{equation*}
        D=
        \begin{mat}
          d_1  &0   &0   &\ldots    \\
          0    &d_2 &0   &          \\
          0    &0   &d_3            \\
               &    &    &\ddots    \\
               &    &    &      &d_n   
        \end{mat}
      \end{equation*}
      若 $i\neq j$，那么 $i,j$~子式是一个 $\nbyn{(n-1)}$ 矩阵，
      其中只有 $n-2$ 个非零元，因为我们同时删去了 $d_i$ 与 $d_j$。
      于是，该子式至少有一行或一列全为零，所以代数余子式 $D_{i,j}$ 为零。
      若 $i=j$，那么该子式是对角元为
      $d_1$, \ldots, $d_{i-1}$, $d_{i+1}$, \ldots, $d_n$ 的对角矩阵。
      它的行列式显然等于 $(-1)^{i+j}=(-1)^{2i}=1$ 
      乘以这些元的乘积。
      \begin{equation*}
        \adj(D)
        =
        \begin{mat}
          d_2\cdots d_n    &0                 &      &0    \\
          0                &d_1d_3\cdots d_n  &      &0    \\
                           &                  &\ddots      \\
                           &                  &       &d_1\cdots d_{n-1} 
        \end{mat}
      \end{equation*}

      顺便说一下，\nearbytheorem{th:MatTimesAdjEqDiagDets} 提供了
      一种更简洁的方法来推导这个结论。
    \end{answer}
  \recommended \item
    证明伴随矩阵的转置等于转置的伴随矩阵。
    \begin{answer}
      只需注意，若 $S=\trans{T}$，则代数余子式
      $S_{j,i}$ 等于代数余子式 $T_{i,j}$，这是因为 $(-1)^{j+i}=(-1)^{i+j}$，
      并且两个子式互为转置（而转置的行列式等于原矩阵的行列式）。
    \end{answer}
  \item 
    证明或证伪：\( \adj(\adj(T))=T \)。
    \begin{answer}
      该命题不成立；这里有一个例子。
      \begin{equation*}
         T=\begin{mat}[r]
           1  &2   &3  \\
           4  &5   &6  \\
           7  &8   &9
         \end{mat}
         \qquad
         \adj(T)=\begin{mat}[r]
          -3  &6   &-3 \\
           6  &-12 &6  \\
          -3  &6   &-3
         \end{mat}
         \qquad
         \adj(\adj(T))=\begin{mat}[r]
           0  &0   &0  \\
           0  &0   &0  \\
           0  &0   &0
         \end{mat}
      \end{equation*}
    \end{answer}
\item
    若方阵的每个 \( i,j \) 元在对角线上方的部分为零，即当 \( i>j \) 时为零，
    则称它是\definend{上三角矩阵}\index{matrix!triangular}。
    \begin{exparts}
      \partsitem
        上三角矩阵的伴随矩阵必定是上三角的吗？
        是下三角的吗？
      \partsitem 证明：若上三角矩阵存在逆矩阵，
        则其逆矩阵是上三角的。
    \end{exparts}
    \begin{answer}
      \begin{exparts}
        \partsitem 例子 
          \begin{equation*}
             M=
            \begin{mat}[r]
              1  &2  &3  \\
              0  &4  &5  \\
              0  &0  &6
            \end{mat}
          \end{equation*}
          提示了正确的答案。
          \begin{align*}
            \adj(M)=
            \begin{mat}
             M_{1,1}  &M_{2,1}  &M_{3,1}  \\
             M_{1,2}  &M_{2,2}  &M_{3,2}  \\
             M_{1,3}  &M_{2,3}  &M_{3,3}  
            \end{mat}
            &=\begin{mat}
              \begin{vmat}
                4  &5 \\ 0 &6
              \end{vmat}
              &-\begin{vmat}
                2  &3  \\  0  &6
              \end{vmat}
              &\begin{vmat}
                2  &3  \\  4  &5 
              \end{vmat}          \\[1.5ex]
              -\begin{vmat}
                0  &5  \\  0  &6
              \end{vmat}
              &\begin{vmat}
                1  &3  \\  0  &6
              \end{vmat}
              &-\begin{vmat}
                1  &3  \\  0  &5  
              \end{vmat}           \\[1.5ex]
              \begin{vmat}
                0  &4  \\  0  &0
              \end{vmat}
              &-\begin{vmat}
                1  &2  \\  0  &0
              \end{vmat}
              &\begin{vmat}
                1  &2  \\  0  &4
              \end{vmat}
            \end{mat}                                \\
            &=
            \begin{mat}[r]
              24  &-12 &-2 \\
               0  &6   &-5  \\
               0  &0   &4
            \end{mat}
          \end{align*}
          结果确实是上三角的。

          这项检验是细致的，但并不难。
          伴随矩阵上三角部分的元是
          $M_{a,b}$，其中 $a>b$。
          我们需要验证：若 $a>b$，则代数余子式 $M_{a,b}$ 为零。
          当 $a>b$ 时，把 $M$ 的第~$a$ 行与第~$b$ 列
          \begin{equation*}
            \begin{mat}
              m_{1,1} &\ldots   &m_{1,b}  &\ldots       &        \\
              m_{2,1} &\ldots   &m_{2,b}  &       &        \\
              \vdots  &         &\vdots   &       &        \\
              m_{a,1} &\ldots   &m_{a,b}  &\ldots &m_{a,n} \\
              \vdots        &         &\vdots   &       &        \\
                      &         &m_{n,b}  &       &                   
            \end{mat}
          \end{equation*}
          删去后，得到一个上三角的子式，
          这是因为该子式的第~$i,j$ 元，要么是 $M$ 的第~$i,j$ 元
          （当 $a>i$ 且~$b>j$ 时发生这种情形；
          此时 $i<j$ 蕴含该元为零），
          要么是 $M$ 的第~$i,j+1$ 元（当 $i<a$ 且
          $j>b$ 时发生这种情形；此时 $i<j$ 蕴含 $i<j+1$，从而
          该元为零），要么是 $M$ 的第~$i+1,j+1$ 元
          （最后这种情形在 $i>a$ 且~$j>b$ 时发生；显然这里
          $i<j$ 蕴含 $i+1<j+1$，所以该元为零）。 
          于是该子式的行列式就是沿对角线的乘积。
          注意，$M$ 的 $a-1,a$ 元就是该子式的
          $a-1,a-1$~元（由于关系 $a>b$ 是严格的，所以它
          不会被删去）。
          但这个元为零，因为 $M$ 是上三角的且
          $a-1<a$。
          因此该代数余子式为零，伴随矩阵是上三角的。
          （下三角的情形类似。）  
        \partsitem 利用\nearbytheorem{th:MatTimesAdjEqDiagDets}，
          这可由前一小问立即得出。
      \end{exparts}
    \end{answer}
  \item 
    \label{exer:LaplaceProof}
    \textit{本题需要选读小节「行列式的存在性」中的内容。}
    用置换展开证明\nearbytheorem{th:LaPlaceExp}。
    \begin{answer}
      我们将证明每个行列式都可以沿第~\( i \) 行展开。
      对第~$j$ 列的论证是类似的。

      置换展开中的每一项恰好含有来自每一行的一个元。
      如\nearbyexample{ex:ExpThreeFirstRow}那样，
      把每个第~$i$ 行的元提出来，得到
      \( \deter{T}
           =t_{i,1}\cdot\hat{T}_{i,1}+\dots+t_{i,n}\cdot\hat{T}_{i,n} \)，
      其中每个 \( \hat{T}_{i,j} \) 是一些不含第 \( i \) 行任何元素的项之和。
      我们将证明 \( \hat{T}_{i,j} \) 就是 \( i,j \) 代数余子式。

      先考虑 \( i,j=n,n \) 的情形：
      \begin{equation*}
        t_{n,n}\cdot\hat{T}_{n,n}
        =t_{n,n}\cdot
            \sum_{\phi}t_{1,\phi(1)}t_{2,\phi(2)}\dots\,t_{n-1,\phi(n-1)}
                      \sgn(\phi)
      \end{equation*}
      其中求和取遍所有满足
      \( \phi(n)=n \) 的 \( n \)-置换 \( \phi \)。
      为了证明
      \( \hat{T}_{i,j} \) 是子式 \( T_{i,j} \)，我们只需证明：
      若 \( \phi \) 是满足 
      \( \phi(n)=n \) 的 \( n \)-置换，
      且 \( \sigma \) 是满足
      \( \sigma(1)=\phi(1) \), \ldots, \( \sigma(n-1)=\phi(n-1) \)
      的 \( n-1 \)-置换，
      则 \( \sgn(\sigma)=\sgn(\phi) \)。
      但这是成立的，因为 $\phi$ 与 $\sigma$ 
      有相同数目的逆序。

      回到一般的 \( i,j \) 情形。
      对换相邻的行，直到第 \( i \) 行变为最后一行；再对换相邻的列，
      直到第 \( j \) 列变为最后一列。
      注意，第 \( i,j \) 个子式的行列式不受这些相邻
      对换的影响，因为逆序被保持（由于该子式
      略去了第 \( i \) 行与第 \( j \) 列）。
      另一方面，\( \deter{T} \) 与 \( \hat{T}_{i,j} \) 的符号
      改变了 \( n-i \) 加 \( n-j \) 次。
      于是 \( \hat{T}_{i,j}=(-1)^{n-i+n-j}\deter{T_{i,j}}
                =(-1)^{i+j}\deter{T_{i,j}} \)。  
    \end{answer}
  \item 
    利用 Laplace 展开以及对矩阵阶数的归纳，
    证明矩阵的行列式等于其转置的行列式。
    \begin{answer}
      对 \( \nbyn{1} \) 的基准情形，这是显然的。

      归纳的情形：
      假设对所有 \( \nbyn{1} \), \ldots, \( \nbyn{(n-1)} \)
      矩阵，矩阵的行列式都等于其转置的行列式。
      沿第~\( i \) 行展开给出
      \( \deter{T}=t_{i,1}T_{i,1}+\dots\,+t_{i,n}T_{i,n} \)，
      而沿第~\( i \) 列展开给出
      \( \deter{\trans{T}}=t_{1,i}(\trans{T})_{1,i}
           +\dots+t_{n,i}(\trans{T})_{n,i} \)。
      由于 \( (-1)^{i+j}=(-1)^{j+i} \)，两个和式中的符号是相同的。
      由于 \( \trans{T} \) 的 \( j,i \) 子式是 \( T \) 的 \( i,j \) 子式的转置，
      归纳假设给出 \( \deter{(\trans{T})_{i,j}}=\deter{T_{i,j}} \)。  
    \end{answer}
  \puzzle \item 
    证明
    \begin{equation*}
      F_n=
      \begin{vmat}
        1  &-1  &1  &-1  &1  &-1  &\ldots  \\
        1  &1   &0  &1   &0  &1   &\ldots  \\
        0  &1   &1  &0   &1  &0   &\ldots  \\
        0  &0   &1  &1   &0  &1   &\ldots  \\
        .  &.   &.  &.   &.  &.   &\ldots
      \end{vmat}
    \end{equation*}
    其中 \( F_n \) 是
    \( 1,1,2,3,5,\dots,x,y,x+y,\ldots\, \)——Fibonacci 数列——的第 \( n \) 项，
    且该行列式的阶为 \( n-1 \)。
    \cite{Monthly49p409}
    \begin{answer}
      \answerasgiven %
      把上面的行列式记作 \( D_n \)，可以看到
      \( D_2=1 \)，\( D_3=2 \)。
      剩下要证明的是 \( D_n=D_{n-1}+D_{n-2},\; n\geq 4 \)。
      在 \( D_n \) 中，从第 \( (n-1) \) 列减去第 \( (n-3) \) 列，
      从第 \( (n-2) \) 列减去第 \( (n-4) \) 列，\ldots，从第三列减去
      第一列，得到
      \begin{equation*}
        F_n=
        \begin{vmat}
          1  &-1  &0  &0   &0  &0   &\ldots  \\
          1  &1   &-1 &0   &0  &0   &\ldots  \\
          0  &1   &1  &-1  &0  &0   &\ldots  \\
          0  &0   &1  &1   &-1 &0   &\ldots  \\
          .  &.   &.  &.   &.  &.   &\ldots
        \end{vmat}.
      \end{equation*}
      按第一行展开这个行列式，即得所要的关系式。  
    \end{answer}
%  \recommended \item 
%    Prove that a square matrix is singular if and only if its adjoint is
%    singular.
%  \item 
%    Prove that if \( T \) is an \( \nbyn{n} \) matrix then
%    \( \det(\adj(T))=(\det(T))^{n-1} \).
%    \begin{answer}
%      By \nearbytheorem{th:MatTimesAdjEqDiagDets}, 
%      $\deter{\adj(T)\cdot T}=\deter{T}^n$ because 
%      $\deter{T}\cdot I$ is a diagonal matrix.
%      Then $\deter{T}^n=\deter{\adj(T)\cdot T}
%                         =\deter{\adj(T)}\cdot \deter{T}$.
%    \end{answer}
\end{exercises}
\index{Laplace determinant expansion|)}
